\section{Introduction}\label{sec:introduction}

For \(n\ge1\), let \(\Omega(n)\) be the number of prime factors of \(n\),
counted with multiplicity, and let \(\lambda(n)=(-1)^{\Omega(n)}\).
We study ordinary two-point averages: each integer up to the given
cutoff has weight one. Thus the correlation measures whether the two prime-factor counts
have the same parity more often than different parities on the full
initial interval.
The distinction from a logarithmic average is the distinction between
\[
 \frac1X\sum_{n\le X}a(n)
 \qquad\text{and}\qquad
 \frac1{\log X}\sum_{n\le X}\frac{a(n)}n.
\]
Logarithmic weighting combines information from many multiplicative
scales; an ordinary estimate must control the terminal scale itself.
Our quantitative conclusion concerns every fixed pair of
nonproportional affine forms.

\begin{theorem}[A logarithmic saving for affine Liouville correlations]
\label{thm:q-affine}
There is an absolute constant \(c>0\) such that, for every fixed choice
of integers \(a_1,a_2\ge1\) and \(b_1,b_2\ge0\) with
\(a_1b_2-a_2b_1\ne0\), there is a constant
\(C_{a_1,a_2,b_1,b_2}\) for which
\[
 \left|\sum_{1\le n\le X}
       \lambda(a_1n+b_1)\lambda(a_2n+b_2)\right|
 \le C_{a_1,a_2,b_1,b_2}\frac{X}{(\log X)^c}
 \qquad(X\ge3).
\]
The cutoff \(X\) may be any real number. No coprimality condition is
imposed on the coefficients.
\end{theorem}

For $a_1=a_2=1$ and $b_1=0$, $b_2=h>0$, this resolves the
ordinary two-point Chowla conjecture, with a logarithmic saving.
The exponent is independent of the forms. The implied constants need
not be effective, and no uniformity for coefficients growing with \(X\)
is asserted. The proof first obtains the same saving for
\(\lambda(n)\lambda(n+h)\) in each fixed residue class modulo each fixed
positive integer, including nonunit and zero classes. Complete
multiplicativity then gives the affine conclusion.

We also prove a qualitative statement for general multiplicative
functions. Let \(\mathbb D=\{z\in\C:|z|\le1\}\). A function
\(f:\N\to\mathbb D\) is multiplicative when
\(f(mn)=f(m)f(n)\) for coprime positive integers \(m,n\).
For a Dirichlet character \(\chi\), a real number \(t\), and \(X\ge2\),
define the nonnegative distance \(D\) by
\begin{equation}\label{eq:distance}
 D(f,\chi n^{it};X)^2
 =\sum_{p\le X}\frac{1-\operatorname{Re}
       (f(p)\overline{\chi(p)p^{it}})}p.
\end{equation}
All sums indexed by \(p\) are over primes. We call \(f\)
\emph{uniformly nonpretentious} if, for every fixed Dirichlet character,
\begin{equation}\label{eq:unp}
 \inf_{|t|\le N}D(f,\chi n^{it};N)\longrightarrow\infty
 \qquad(N\longrightarrow\infty).
\end{equation}
The prime cutoff and the height bound are the same \(N\).

\begin{theorem}[Binary corrected Elliott]\label{thm:elliott}
Let \(f_1,f_2:\N\to\mathbb D\) be multiplicative, and suppose that
at least one satisfies \eqref{eq:unp}. For every fixed pair of distinct
nonnegative integers \(h_1,h_2\),
\[
 \frac1N\sum_{n=1}^N f_1(n+h_1)f_2(n+h_2)\longrightarrow0
 \qquad(N\longrightarrow\infty,\ N\in\N).
\]
\end{theorem}

This proves the binary corrected Elliott conjecture for ordinary
averages. The functions may be complex, need not have modulus one,
and need not be completely multiplicative. There is no conjugation
in the displayed product; the conjugated version follows because
\eqref{eq:unp} is preserved by conjugation. We make no quantitative
rate claim for this class of functions.

\begin{corollary}[Affine correlations]\label{cor:affine}
Let \(f_1,f_2:\N\to\mathbb D\) be multiplicative, with at least one
satisfying \eqref{eq:unp}. For all fixed integers \(a_1,a_2\ge1\) and
\(b_1,b_2\ge0\) with \(a_1b_2-a_2b_1\ne0\),
\[
 \frac1N\sum_{n=1}^N f_1(a_1n+b_1)f_2(a_2n+b_2)
 \longrightarrow0.
\]
\end{corollary}

Every fixed progression restriction is also permitted in
\Cref{thm:elliott}; \Cref{prop:affine-progression} gives the precise
statement through all real cutoffs and for every residue class.
These conclusions use the hypothesis on an original factor \(f_j\).
The finite local changes needed for dilations and residue restrictions
preserve that hypothesis. In particular, Liouville satisfies the exact
condition \eqref{eq:unp}; a quantitative distance bound is proved in
\Cref{lem:affine-liouville-unp}.

\subsection{Historical context and prior work}

Chowla's conjecture predicts vanishing Liouville correlations at every
fixed collection of distinct shifts \cite{Chowla65,MRT15}. Its binary
case asks whether the parities of the numbers of prime factors of
\(n\) and \(n+h\) become uncorrelated for each fixed \(h\ge1\).
Partial summation transfers ordinary cancellation to logarithmic
cancellation. An ordinary estimate at every cutoff requires more
than this implication.

Elliott extended the correlation problem to bounded multiplicative
functions \cite{Elliott92}. The obstruction presented by the twists
\(n^{it}\) is already central to Hal\'asz's theory of one-point mean
values \cite{Halasz71}; see also the treatment of Granville, Harper,
and Soundararajan \cite{GHS19}. The distance in \eqref{eq:distance}
belongs to the pretentious approach developed by Granville and
Soundararajan \cite[Section~3]{GS07}. Matom\"aki, Radziwi\l\l, and Tao
showed why Elliott's original condition needs correction for complex
functions: a function can imitate different twists at successive
scales without imitating any one fixed twist globally. Their
corrected condition is equivalent to \eqref{eq:unp}; see
\cite[Section~1.1 and Appendix~B]{MRT15}.

Matom\"aki and Radziwi\l\l\ proved that most short means of a real
bounded multiplicative function approximate its long mean
\cite[Theorem~1]{MR16}. They also obtained a nontrivial bound for each
fixed-shift Liouville correlation: the normalized absolute value is
at most \(1-\delta(h)\), for some \(\delta(h)>0\), at all sufficiently
large cutoffs \cite[Corollary~2]{MR16}. Matom\"aki, Radziwi\l\l, and
Tao then established cancellation averaged over arbitrarily slowly
growing windows of shifts. Their exponential-sum extension to
complex multiplicative functions supplies the short-interval
estimates used here \cite[Theorems~1.1, 1.3, and~1.7]{MRT15}.

Tao proved the logarithmically averaged binary corrected Elliott
theorem for nonproportional affine forms
\cite[Theorem~1.3 and Corollary~1.5]{Tao16}.
For the logarithmically weighted Liouville sum, Helfgott and
Radziwi\l\l\ obtained
\[
 \left|\sum_{n\le x}\frac{\lambda(n)\lambda(n+1)}n\right|
 \ll \frac{\log x}{\sqrt{\log\log x}}
\]
by divisibility-graph expansion \cite[Corollary~1.5]{HR21}.
Pilatte improved this bound to \((\log x)^{1-c_0}\) for an absolute
\(c_0>0\) \cite[Theorem~1.1]{Pilatte26}.
These estimates retain the logarithmic weight.

The graph method of Helfgott and Radziwi\l\l\ uses high traces,
recurrence constraints, and cancellation at primes occurring only
once \cite[Sections~2--7]{HR21}. They also proposed composite steps
and nonbacktracking operators \cite[Sections~9.1--9.2]{HR21}.
Pilatte develops these directions through products of centered prime
factors, prohibited sequences, triangular arithmetic constraints,
and a rank-truncated intersection sieve
\cite[Sections~2--13 and Appendix~B]{Pilatte26}.
Our two graph arguments adapt these constructions to different
weights and orders of limits. Their trace estimates, local
expansions, and transfers to ordinary averages are proved below.

For ordinary averages,
Tao and Ter\"av\"ainen obtained binary cancellation outside
a set of scales of logarithmic density zero, under the same
condition \eqref{eq:unp} on one factor
\cite[Corollary~1.13]{TT19}. Klurman, Mangerel, and Ter\"av\"ainen
relaxed the hypothesis to nonpretentiousness against each fixed
\(\chi(n)n^{it}\),
obtaining cancellation along a set of scales of full upper logarithmic
density \cite[Theorem~1.2]{KMT23}. The later quantitative progression
estimates of Tao and Ter\"av\"ainen
give logarithmic savings for affine Liouville correlations,
with coefficients allowed to grow slowly, outside a quantitatively
controlled exceptional set of scales
\cite[Theorem~3.1 and Remarks~3.2]{TT26}. The results here concern
every cutoff, with the affine forms fixed.

The quantitative route also uses bounded independence for Boolean
circuits. This subject developed from the work of Linial and Nisan
\cite{LinialNisan90}, through Bazzi's depth-two theorem and
Razborov's simpler proof \cite{Bazzi09,Razborov08}, to Braverman's
theorem for every fixed depth \cite{Braverman10}. In applying that theorem, the
encoded residue law is first corrected to exact bounded independence.
The local Fourier calculation is a form of the almost-to-exact
independence construction of Alon, Goldreich, and Mansour
\cite[Theorem~2.1 and Section~3.1]{AGM03}.

\subsection{Main ideas}

Both arguments begin with multiplicative dilations. For a fixed
shift $h\ge1$, comparison of the correlations at $n,n+h$ and at
$dn,d(n+h)$ leads to a graph on the integers with edge displacement
$hd$. Short-interval Fourier estimates control the change in the
correlation when selected divisibility indicators are replaced by
signed expressions. Estimates for closed walks then bound the resulting
graph operator. The two proofs use different vertex weights, and their
cancellation mechanisms must be matched to those weights.

In the quantitative argument, a step contains one prime from each of
$J$ disjoint supplies, each with reciprocal mass comparable to a large
fixed constant $W$. These primes contribute the factors
$\1_{p\mid n}-1/p$, which have mean zero under uniform residues.
An additional squarefree padding divisor uses a disjoint set of primes.
The centered supplies provide reciprocal mass at least $W^J$, while
their contribution to the operator bound is only $(C\sqrt W)^J$, with
an absolute $C$. Padding controls the concentration
of divisor weight in short multiplicative intervals. Writing $L$ for
a small fixed power of $\log X$, the edge normalization gives a factor
$1/L$ in the estimate for each interval, compensating for the factor
$L$ in the number of intervals. The remaining losses are exponential
in $J$ with absolute bases. Choosing $W$ large enough therefore gives
exponential decay in $J$; since $J$ grows proportionally to
$\log\log X$, this is a power-of-logarithm saving.

Two issues underlie this comparison. First, the prime supplies grow
with $X$, so averaging over their full joint period would not give
a useful finite-scale estimate. We instead compare the residue tests
used in the proof with independent residues, correcting their encoded
bit law to exact bounded independence before applying Braverman's
theorem \cite{Braverman10}. Second, repeated primes constrain the
closed walks. Many independent arithmetic relations give a direct
saving. When few remain, a forest encodes the possible patterns of
repetition. Its pattern count is uniform over all padding values,
allowing the padding weights to be summed afterward in one common
residue environment. These steps preserve the absolute constants
needed for the choice of $W$.

For general multiplicative functions, the short-interval input
contains the nonpretentiousness distance, whose divergence has no
prescribed rate. We therefore fix the graph before taking the long
average. At an auxiliary scale $B$, squarefree divisors are formed from
two bands of primes at most $e^B$. Primes in the higher band receive a
fixed weight greater than one, and those in the lower band receive
weight one; the weight of a divisor is the product of its prime weights.
After dividing the ordinary divisor average by the total reciprocal
weight of these divisors, a biased correlation subsequence produces a
raw average of size at least a constant times $B^{-1}$ in one short
multiplicative interval. The analytic centering error is $o(B^{-1})$,
and the graph estimate gives $O(B^{-1-c_* /2})$ for a fixed absolute
$c_*>0$. Taking the long limit with $B$ fixed, and then increasing $B$,
contradicts the bias.

The higher prime band, called the core band, supplies decay through
restrictions on divisor size and on the number of core primes at a
vertex. The lower, center band provides the additional power saving
needed to beat the $B^{-1}$ scale. Its signed factor is
$\1_{p\mid n}-\theta/p$, with
$\theta$ chosen to match the vertex normalization; it is not itself
mean zero under uniform residues. For suitable primes appearing on
only one edge, the normalization factors at the two endpoints in the
closed-walk product make the leading divisibility contribution match
the subtracted constant term.
To retain this cancellation,
we organize a walk by the tree of its first visits and expand the
vertex-deletion conditions while leaving these center-prime residues
unconditioned. Their signed averages are taken before absolute values.
This yields the qualitative theorem without a Liouville-specific rate.

\subsection{Organization}

Part~I proves the quantitative progression estimate and then
\Cref{thm:q-affine}.
\Cref{q:setup,q:finite-law-section} construct its prime supplies and
establish the finite residue comparison; \Cref{q:analytic} bounds
centering and deletion errors. \Cref{q:paths-trace} introduces the
graph moment, and \Cref{q:rank-witness-section,q:forest-section} prove
it through arithmetic constraints and pattern counting.
\Cref{q:sec-spectral} returns to ordinary correlations, while
\Cref{q:affine-transfer} gives the affine deduction.

Part~II proves \Cref{thm:elliott} with new parameters and prime sets.
\Cref{sec:graph-setup,sec:analytic-transfer,sec:analytic-centering}
state the graph estimate and prove the analytic reduction to it.
\Cref{sec:paths,sec:thinning} establish the geometry of retained walks;
\Cref{sec:cylinder-sieve,sec:trace-count} control their conditional
averages and sum the trace. \Cref{sec:localization} transfers that
estimate to long integer intervals and completes the theorem.
Part~III uses finite local expansions to prove the progression and
affine consequences for general functions, and verifies their
Liouville and M\"obius specializations.
